% Compléments Bloc 3 — zooms exploitation / industrialisation

\begin{frame}{C3.1 + C3.2 — Une règle simple : pas de build sans validation}
\context{B3}{C3.1 · C3.2 · piv2glz · Zoom preuve}
\begin{center}
\begin{tikzpicture}[node distance=0.9cm,every node/.style={font=\small,align=center}]
  \node[draw=Line,rounded corners,fill=SoftBlue,minimum width=2.5cm,minimum height=0.9cm] (commit) {Commit};
  \node[draw=Line,rounded corners,fill=SoftGreen,minimum width=2.5cm,minimum height=0.9cm,right=of commit] (tests) {Pytest};
  \node[draw=Line,rounded corners,fill=SoftOrange,minimum width=2.5cm,minimum height=0.9cm,right=of tests] (build) {PyInstaller};
  \node[draw=Line,rounded corners,fill=SoftPurple,minimum width=2.5cm,minimum height=0.9cm,right=of build] (art) {Artefacts};
  \draw[->,thick,Blue] (commit)--(tests); \draw[->,thick,Blue] (tests)--(build); \draw[->,thick,Blue] (build)--(art);
\end{tikzpicture}
\end{center}
\vspace{0.8em}
\begin{tightitems}
  \item La CI ne sert pas seulement à automatiser : elle rend la règle de validation explicite.
  \item Le livrable produit correspond à un état du dépôt ayant franchi le même contrôle reproductible.
  \item Une automatisation sans test peut simplement automatiser une régression.
\end{tightitems}
\end{frame}

\begin{frame}{C3.3 — Surveiller une chaîne, c'est aussi redécouvrir ses règles métier}
\context{B3}{C3.3 · FCU / Ediacara · Zoom preuve}
\begin{columns}[T]
\begin{column}{0.46\linewidth}
\begin{block}{Boucle d'exploitation}
\begin{enumerate}
  \item Observer les jobs et les logs.
  \item Qualifier l'anomalie.
  \item Remonter jusqu'à la donnée ou au composant source.
  \item Corriger puis vérifier en production.
\end{enumerate}
\end{block}
\end{column}
\begin{column}{0.50\linewidth}
\begin{block}{Cas rencontrés}
\begin{tightitems}
  \item Champ PDF \texttt{/UserUnit}.
  \item Désynchronisation possible entre \texttt{.lis} et ShomProduit.
  \item Règles de synchronisation implicites révélées par le FCU.
\end{tightitems}
\end{block}
\end{column}
\end{columns}
\begin{block}{Message jury}
Une supervision utile relie l'état technique du job au processus métier qui produit et consomme la donnée.
\end{block}
\end{frame}

\begin{frame}{C3.6 — Plusieurs documents pour plusieurs usages}
\context{B3}{C3.6 · FCU / Ediacara · Zoom preuve}
\scriptsize
\begin{tabularx}{\linewidth}{>{\bfseries}p{3.0cm}X X}
\toprule
Support & Fonction & Besoin couvert \\
\midrule
Diagramme global & visualiser les flux et dépendances & localiser rapidement une panne \\
MkDocs & décrire jobs, entrées, sorties, commandes & maintenir ou modifier un traitement \\
Matrice PROD / RE7 & comparer les environnements & sécuriser la recette \\
Procédure de MEP & exécuter et revenir en arrière & maîtriser une intervention sensible \\
Document ArchFCU & conserver le raisonnement initial & expliquer choix et alternatives \\
\bottomrule
\end{tabularx}
\vspace{0.7em}
\begin{block}{Point de recul}
J'ai d'abord documenté pour comprendre ; la transmission à l'équipe est devenue ensuite un résultat concret de ce travail.
\end{block}
\end{frame}

\begin{frame}{Numérique responsable — L'efficience avant le discours}
\context{B3}{Compétence transversale · Zoom preuve}
\begin{columns}[T]
\begin{column}{0.48\linewidth}
\begin{block}{Décisions concrètes}
\begin{tightitems}
  \item Éviter les recopies systématiques de raster volumineux.
  \item Comparer taille, date ou empreinte avant transfert.
  \item Utiliser des tampons locaux avant livraison réseau.
  \item Réutiliser du matériel existant pour le homelab.
\end{tightitems}
\end{block}
\end{column}
\begin{column}{0.48\linewidth}
\begin{block}{Ce que je ne prétends pas}
\begin{tightitems}
  \item Pas de bilan carbone complet.
  \item Pas d'analyse de cycle de vie.
  \item Pas de mesure énergétique instrumentée de la chaîne.
\end{tightitems}
\end{block}
\end{column}
\end{columns}
\begin{block}{Positionnement}
Je présente ces choix comme des décisions d'efficience technique mesurées dans leur portée, pas comme une preuve environnementale globale.
\end{block}
\end{frame}
